\section{Introduction}
\label{sec:introduction}

\subsection{The Over-Optimising Objection}

\cs{}~\citep{b16paper} solves a genuine problem: a fixed single seed for
the METIS heuristic may land in a local optimum, leaving a more compact map
undiscovered.
By sweeping $T = 600$ seeds and returning the global minimum, \cs{}
guarantees that the plan produced is the best available within the bisection
family.

This guarantee is also a potential vulnerability.
A sophisticated opponent could raise what we call the \emph{over-optimising
objection}: the algorithm is specifically designed to find the most extreme
member of its plan family.
In the compactness dimension, this extremum sits at the 0.1--0.7th percentile
of the full \recom{} ensemble for most states (G.1~\citep{g1paper}).
The plan is, by construction, more compact than 99\%+ of all valid
redistricting plans.
No legislature drawing districts by hand would produce anything close to
this level of compactness.
Could this extremism itself be a form of manipulation---a gerrymander of
compactness rather than partisanship?

The over-optimising objection is legally novel.
Traditional gerrymandering doctrine targets partisan or racial distortion.
\rucho{}~\citep{rucho2019} forecloses federal judicial review of partisan
claims; \wesberry{}~\citep{wesberry1964} and \karcher{}~\citep{karcher1983}
establish the population-equality framework.
No existing doctrine directly addresses whether maximising compactness is
itself a legally cognisable harm.
Nevertheless, the objection is worth anticipating.

\subsection{The PercentileSweep Response}

\ps{} dissolves the over-optimising objection by making the compactness
percentile explicit.
Instead of always returning the minimum-edge-cut plan, \ps{}$(p, T)$ returns
the plan at percentile $p$ of the $T$-plan bisection distribution.
The percentile $p$ is a statutory parameter, visible in the legislation,
subject to legislative debate, and auditable by any party.

This has three immediate legal advantages.
First, the choice of percentile is \emph{transparent}: opponents and
supporters alike can see exactly what posture the statute adopts.
Second, the choice is \emph{principled}: each value of $p$ corresponds to
a recognisable legal doctrine (Section~\ref{sec:legal-posture}).
Third, the choice is \emph{auditable}: given the statutory $p$, any party
can reproduce the plan by running the same $T$ seeds and selecting rank
$\floor{p \times T}$.

\subsection{The Insensitivity Finding}

Critically, the choice of $p$ \emph{does not matter much} for partisan
outcomes.
Our empirical analysis across six key states (NC, WI, GA, PA, TX, CA) shows
that partisan seat counts vary by at most 0.5 seats across the full range
$p \in \{0.0, 0.25, 0.5, 0.75, 1.0\}$ with $T = 101$ plans.

This finding extends B.17's parameter insensitivity result~\citep{b17paper}
to the search percentile dimension.
B.17 showed that partisan outcomes are insensitive to METIS parameters,
convergence threshold $T$, and edge-weight settings within a fixed algorithm
family.
H.0 shows they are equally insensitive to which compactness level within
the bisection family is selected.
The underlying mechanism is the same: geographic structure, not search
strategy, determines partisan outcomes (B.7~\citep{b7paper}).

The insensitivity finding is both a scientific result and a legal asset.
It means that the statutory choice of percentile can be made on legal and
normative grounds without partisan consequences.
The DIA can adopt $p = 0.0$ because it produces the most legally defensible
compactness posture, not because it manipulates partisan outcomes.

\subsection{Paper Structure}

Section~\ref{sec:related-work} situates \ps{} within the B-series
and G-series literature.
Section~\ref{sec:algorithm} defines \ps{} formally.
Section~\ref{sec:results} presents empirical results for six states.
Section~\ref{sec:legal-posture} maps percentile choices to legal doctrine.
Section~\ref{sec:conclusion} concludes with a statutory recommendation.
